% ===================================================================
%  કોષ્ટક નં. ૧૨ — સ્ટેશન. --- રેલવે. --- જહાજ.
%
%  Traduction, faite en 2026, en gujarati d'aujourd'hui, sur l'ido de
%  texto/io. Regles de la colonne : voir 10-kostak-01.tex. Une seule
%  numerotation. Le « Noto. » d'ouverture fait partie du bloc apar, il
%  n'a pas de cle a lui ; la note de bas de page, elle, en a une —
%  t12-noto-1, posee entre le bloc 14 et le bloc 15, comme dans l'ido.
%
%  LE CHEMIN DE FER EST LA SEULE CHOSE MODERNE QUE L'INDE AIT RECUE EN
%  MEME TEMPS QUE L'EUROPE, ET LA LANGUE LE MONTRE. Le premier train
%  indien a roule en 1853, vingt-huit ans apres le premier anglais et
%  bien avant que la France du livret n'ait son reseau. Le vocabulaire
%  n'est donc ni un vieux fonds ni un emprunt tardif : c'est un
%  partage fait a la meme date. Et le partage est net. Les trois
%  pieces qu'on peut voir et toucher ont recu un mot d'ici, pris a
%  quelque chose qui existait avant : « ગાડી » le train, qui etait la
%  charrette ; « ડબો » le wagon, qui etait la boite ; « પાટા » les
%  rails, qui etaient les planches. Ce qu'on ne peut qu'apprendre est
%  reste anglais : સિગ્નલ, પ્લૅટફૉર્મ, ટિકિટ, ગાર્ડ, એન્જિન.
%  UNE LANGUE REBAPTISE CE QU'ELLE VOIT ET EMPRUNTE CE QU'ON LUI
%  EXPLIQUE.
%
%  ET UN SEUL MOT DE CE TABLEAU A FAIT LE VOYAGE EN SENS INVERSE :
%  « કુલી » (16). Le porteur de gare est le seul terme de tout le
%  releve qui soit parti d'ici pour entrer dans les autres langues —
%  coolie en anglais et en francais, avec la charge que l'on sait. Le
%  livret le grave comme un metier de quai parmi d'autres, et la
%  colonne l'ecrit comme tel, sans note dans le texte : c'est le mot
%  ordinaire, aujourd'hui encore, sur les quais de Ahmedabad. Mais il
%  fallait le dire ici : dix tableaux d'emprunts, et le onzieme rend
%  un pret.
%
%  L'HORLOGE DE LA FRONTIERE (89) AVANCE DE CINQUANTE-CINQ MINUTES, ET
%  C'EST L'ALINEA QUI SE LIT LE PLUS DROLEMENT ICI. Le voyageur de
%  1926 change de pays et doit refaire sa montre. L'Inde a connu la
%  meme chose sans changer de pays : l'heure de Bombay et l'heure de
%  Calcutta ont differe jusqu'en 1906, et Bombay a garde la sienne,
%  a trente-neuf minutes de l'heure indienne, jusqu'en 1955 — les
%  horloges des gares y portaient deux jeux d'aiguilles. Le releve
%  n'invente rien et ne transpose rien : il ecrit la gare de
%  Deutsch-Avricourt et ses cinquante-cinq minutes. Mais le lecteur de
%  cette colonne-ci est le seul du livret a qui la scene ne demande
%  aucun effort d'imagination.
%
%  Le bureau du telegraphe (4) est « તારઘર », la maison du fil, et le
%  telegramme (5) est « તાર », le fil tout court. Le service indien a
%  ferme en 2013, apres cent soixante ans ; le mot, lui, est reste
%  dans « તાર મોકલવો » comme le francais garde « donner un coup de
%  fil » a une epoque ou plus personne ne tourne de cadran. Et le
%  garde-barriere (88) est « ફાટકવાળો », l'homme du ફાટક, le portail
%  du passage a niveau : il y en a des milliers, avec la meme cabane
%  que sur la planche.
% ===================================================================

%%K t12-apar-1 apar
\VUsaut{4.0mm}
\VUtitre{15.0pt}{60}{કોષ્ટક નં. ૧૨}
\VUsaut{2.4mm}
\VUpk{11.4pt}{40}{સ્ટેશન. --- રેલવે. --- જહાજ.}
\VUsaut{3.4mm}
\textsc{નોંધ.} --- દરેક દેશમાં વપરાતાં સમયપત્રક જુદાં જુદાં હોય છે,
એટલે અમે ઉદાહરણ તરીકે \VUgras{ફ્રેન્ચ સમયપત્રકના} કલાક વાપરીએ છીએ.

%%K t12-01-1 p
1. --- હું હમણાં જ જર્મનીમાં થઈને મુસાફરી કરી આવ્યો. નીકળતાં પહેલાં મેં
ગાડીઓ ઊપડવાના કલાક જાણવા \VUgras{સમયપત્રકની ચોપડી} \textsuperscript{(15)}
ખરીદી. સૌથી અનુકૂળ ગાડી પૅરિસથી રાત્રે નવ વાગ્યે ઊપડીને સ્ટ્રાસબુર્ગ
એક વાગ્યે તેર મિનિટે પહોંચે છે એ જોઈ લીધા પછી મેં મારી મુસાફરીની
તૈયારી કરી, અને બીજા દિવસે મેં મને સ્ટેશને પહોંચાડ્યો.

%%K t12-02-1 p
2. --- જ્યારે હું મોટા પ્રસ્થાન-ખંડમાં પેઠો, એક ઘરડા ગામઠી
\VUgras{પાદરીની} \textsuperscript{(2)} પાછળ, જે હાથમાં પોતાની
\VUgras{પ્રવાસની થેલી} \textsuperscript{(3)} લઈને આવ્યા હતા, ત્યારે ઘણા
\VUgras{મુસાફરો} \textsuperscript{(1)} આવી ચૂક્યા હતા.

%%K t12-02-2 p
મારે \VUgras{તાર (ટેલિગ્રામ)} \textsuperscript{(5)} મોકલવો હતો, એટલે મેં
એક \VUgras{ટિકિટ ચકાસનાર} \textsuperscript{(6)} પાસે \VUgras{તારઘર}
\textsuperscript{(4)} બતાવડાવ્યું.

%%K t12-03-1 p
3. --- \VUgras{પ્રતીક્ષાખંડમાં} \textsuperscript{(7)} જતાં પહેલાં હું
પાછા ફરવાની \VUgras{ટિકિટ} \textsuperscript{(9)} લેવા ગયો.

%%K t12-04-1 p
4. --- \VUgras{ટિકિટબારીએ} \textsuperscript{(8)} મારે લાંબી રાહ જોવી ન
પડી, કેમ કે ત્યાં થોડા જ માણસો હતા. એક \VUgras{સિપાહી}
\textsuperscript{(10)}, જે રજા પર જતો હતો, તેણે મને પોતાનો વારો આપવાની
મહેરબાની કરી, એટલે મને \VUgras{સ્ટેશન માસ્તર} \textsuperscript{(13)}
પાસેથી થોડી માહિતી પૂછવાનો પૂરતો સમય મળ્યો, જે દેખરેખના
\VUgras{ઇન્સ્પેક્ટર} \textsuperscript{(12)} સાથે અને ફરજ પરના
\VUgras{હવાલદાર} \textsuperscript{(11)} સાથે વાત કરી રહ્યા હતા.

%%K t12-tit-1 sub
\VUpk{10.2pt}{40}{સામાનની નોંધણી.}
\VUsaut{3.0mm}

%%K t12-05-1 p
5. --- \VUgras{ચોપડીના સ્ટૉલ} \textsuperscript{(14)} પરથી છાપું
ખરીદીને મેં મારો \VUgras{સામાન} \textsuperscript{(17)} જોખાવ્યો, જે એક
\VUgras{કુલી} \textsuperscript{(16)} હમણાં જ \VUgras{સામાનની
હાથલારીમાં} \textsuperscript{(19)} લાવ્યો હતો ; \VUgras{કારકુને}
\textsuperscript{(22)}, જેને \VUgras{વજનકાંટાનું} \textsuperscript{(21)}
કામ સોંપાયું છે, તેણે મને જણાવ્યું કે મારી પેટીનું વજન પાંત્રીસ
કિલોગ્રામ છે ; એટલે પાંચ કિલોગ્રામનું \VUgras{વધારાનું વજન} થયું, જેના
માટે મારે \VUgras{સામાનની બારીએ} \textsuperscript{(23)} વધારાનું ભાડું
ભરવું પડે.

%%K t12-06-1 p
6. --- હું બહુ વહેલો હતો, એટલે મને સ્ટેશનમાં મુસાફરોની અવરજવર જોવાની
ફુરસદ મળી. કેટલાક પોતાનો હલકો સામાન \VUgras{ક્લોકરૂમમાં}
\textsuperscript{(25)} મૂકતા કે પાછો લેતા હતા ; બીજા
\VUgras{સમયપત્રક} \textsuperscript{(26)} જોતા હતા. આવી પહોંચેલા મુસાફરો
\VUgras{બહાર નીકળવાના રસ્તા} \textsuperscript{(32)} તરફ ઉતાવળ કરતા હતા,
હાથમાં \VUgras{બૅગ} \textsuperscript{(24)} લઈને. કેટલાક \VUgras{સામાન
આપવાની જગ્યાએ} \textsuperscript{(27)} રાહ જોતા હતા અને પોતાની
\VUgras{ટિકિટ} \textsuperscript{(29)} બતાવતા હતા, પોતાની ટ્રંક,
\VUgras{પેટીઓ} \textsuperscript{(28)} કે \VUgras{ટોપલા}
\textsuperscript{(30)} પાછા લેવા.

%%K t12-07-1 p
7. --- \VUgras{જકાતના કારકુનો} \textsuperscript{(33)} કાળજીથી પોટલાં
તપાસતા હતા, કંઈ જાહેર કરવા જેવું છે કે નહીં એ જોવા.

%%K t12-08-1 p
8. --- કેટલાક મુસાફરો \VUgras{ઉપાહારગૃહમાં} \textsuperscript{(34)} ગયા,
જ્યાં એક \VUgras{વેઇટર} \textsuperscript{(35)} તેમને પીરસવામાં ઉત્સાહ
બતાવતો હતો. તેમણે ઉતાવળ કરવી પડે, કેમ કે \VUgras{ગાડી}
\textsuperscript{(36)}, જે એક્સપ્રેસ છે, સ્ટેશનમાં લાંબો સમય નહીં રહે.

%%K t12-09-1 p
9. --- પછી હું ઊપડવાના પ્લૅટફૉર્મ પર ગયો અને સીધો \VUgras{ડબો}
\textsuperscript{(37)} શોધ્યો, જેથી રસ્તામાં ડબો બદલવો ન પડે. હું એક
કૉરિડોરવાળી ગાડીમાં ચડ્યો, જેમાં ધૂમ્રપાન કરનારાઓ માટે એક
\VUgras{ખાનું} \textsuperscript{(38)} છે અને « એકલી બહેનો » માટે એક
અલગ ખાનું છે.

%%K t12-tit-2 sub
\VUpk{10.2pt}{40}{રાતનું પ્રસ્થાન.}
\VUsaut{3.0mm}

%%K t12-10-1 p
10. --- \VUgras{પગથિયું} \textsuperscript{(40)} ચડીને,
\VUgras{બારણું} \textsuperscript{(39)} બંધ કરીને, \VUgras{પડદો}
\textsuperscript{(41)} વાળીને અને મારો નાનો સામાન \VUgras{સામાનની
જાળી} \textsuperscript{(43)} પર મૂકીને, હું \VUgras{બાંકડા}
\textsuperscript{(42)} પર બેઠો. \VUgras{દીવાવાળાને}
\textsuperscript{(44)} દીવા સળગાવતો જોઈને મને થયું કે મારે ડબામાં આખી
રાત ગાળવી પડશે, અને મેં \VUgras{ઓશીકાં} \textsuperscript{(45)} ભાડે
આપવાનું કામ જેને સોંપાયું છે તે કારકુનને એક ઓશીકું માગવા બોલાવ્યો.

%%K t12-11-1 p
11. --- ટિકિટ ચેકર ટિકિટ તપાસવા આવ્યો. મેં તેને પૂછ્યું કે અમે હજી કેમ
ઊપડ્યા નથી, કેમ કે સાચો સમય થઈ ગયો છે. તેણે મને જવાબ આપ્યો કે
\VUgras{ગાર્ડ} \textsuperscript{(46)} \VUgras{સામાનના ડબામાં}
\textsuperscript{(47)} સાઇકલો ગોઠવાવવામાં રોકાયેલો છે. વળી, બધાં
\VUgras{પગ ગરમ રાખવાનાં સાધન} \textsuperscript{(48)} હજી બદલાયાં નથી.

%%K t12-12-1 p
12. --- પણ અમે જલદી ઊપડવાના હતા, કેમ કે \VUgras{કોલસો નાખનારે}
\textsuperscript{(50)}, પોતાના \VUgras{ટેન્ડર} \textsuperscript{(49)} પર
ચડીને, \VUgras{એન્જિનની} \textsuperscript{(54)} આગ સતેજ કરવા કોલસો
લીધો, અને \VUgras{ડ્રાઇવર} \textsuperscript{(51)} માત્ર \VUgras{નાયબ
સ્ટેશન માસ્તરના} \textsuperscript{(52)} સંકેતની રાહ જોતો હતો, જે
\VUgras{પ્લૅટફૉર્મ} \textsuperscript{(53)} પર હતા.

%%K t12-13-1 p
13. --- એન્જિનનું \VUgras{બોઇલર} \textsuperscript{(56)} ધીમું ધીમું
ઘૂરકતું હતું ; \VUgras{ભૂંગળું} \textsuperscript{(55)} \VUgras{વરાળ}
\textsuperscript{(60)} સાથે ભળેલા ધુમાડાના ધોધ ઓકતું હતું. કર્કશ
\VUgras{સીટી} \textsuperscript{(59)} ગાજી ઊઠી અને \VUgras{પિસ્ટન}
\textsuperscript{(58)} ચાલવા લાગ્યા, એ ભારે યંત્રનાં \VUgras{પૈડાં}
\textsuperscript{(57)} ધકેલતાં.

%%K t12-tit-3 sub
\VUsaut{3.0mm}
\VUpk{10.2pt}{40}{ઊપડ્યા !}
\VUsaut{3.0mm}

%%K t12-14-1 p
14. --- \VUgras{પાટા} \textsuperscript{(64)} પર દોડતી ગાડીની ઝડપ વધી ;
\VUgras{પ્લૅટફૉર્મ} \textsuperscript{(65)} અદૃશ્ય થયાં ; અમે
\VUgras{સંડાસ} \textsuperscript{(66)} વટાવ્યાં, અને બીજી
\VUgras{લાઇન} \textsuperscript{(63)} પર ઊભા રખાયેલા ડબાઓની હાર પણ :
\VUgras{સૂવાના ડબા} \textsuperscript{(67)}, \VUgras{ભોજનનો ડબો}
\textsuperscript{(68)}, \VUgras{ટપાલનો ડબો} \textsuperscript{(69)}.

%%K t12-noto-1 noto
\VUnotes{143.2mm}{%
\textit{(*) નોંધ.} --- સ્વાભાવિક રીતે આ મુસાફરીનું વર્ણન યુદ્ધ
પહેલાંની સરહદ પ્રમાણે છે.%
}

%%K t12-15-1 p
15. --- પછી \VUgras{માલગોદી} \textsuperscript{(71)} અમારી નજર સામેથી
પસાર થઈ. છાપરાં નીચે કારકુનો \VUgras{માલ} \textsuperscript{(72)} ભરતા
હતા, જે ધીમી ગાડીથી રવાના થવાનો હતો. \VUgras{ટોળીઓ}
\textsuperscript{(76)} \VUgras{પાણીની ટાંકી} \textsuperscript{(73)} પાસે
\VUgras{ઢોરના ડબા} \textsuperscript{(75)} અને \VUgras{સપાટ ડબા}
\textsuperscript{(79)} હલાવીને \VUgras{માલગાડી} \textsuperscript{(80)}
બનાવતી હતી.

%%K t12-16-1 p
16. --- અમે છેલ્લો \VUgras{ફાંટો} \textsuperscript{(78)} વટાવ્યો, જેની
પાસે કાંટાવાળો ઊભો હતો ; અમે \VUgras{સિગ્નલ} \textsuperscript{(86)}
વટાવ્યું અને સ્ટેશન દૂર અદૃશ્ય થઈ ગયું.

%%K t12-17-1 p
17. --- જલદી જ અમે એક \VUgras{લોખંડનો પુલ} \textsuperscript{(74)}
વટાવ્યો, જે \VUgras{નદી} \textsuperscript{(81)} પરથી જાય છે. એ પુલ
વટાવ્યા પછી અમે નદીકાંઠે ચાલ્યા, જેના પર શક્તિશાળી \VUgras{ખેંચનારાં
જહાજ} \textsuperscript{(82)} ભારે \VUgras{માલવાહક હોડીઓ}
\textsuperscript{(83)} ખેંચતાં હતાં. વળી અમે એક \VUgras{મુસાફર જહાજ}
\textsuperscript{(84)} જોયું, જ્યારે તે \VUgras{તરતા ધક્કાને}
\textsuperscript{(85)} અડી રહ્યું હતું.

%%K t12-18-1 p
18. --- કેટલાંક સ્ટેશન પૂરપાટ વટાવ્યા પછી અમે થોડાં \VUgras{બોગદાં}
\textsuperscript{(87)} વટાવ્યાં અને પાછળ છોડ્યાં \VUgras{ફાટકવાળાનાં}
અસંખ્ય \VUgras{નાનાં ઘર} \textsuperscript{(88)}. ગાડીના હલબલાટે
હીંચકતાં હીંચકતાં હું ઘસઘસાટ ઊંઘી ગયો.

%%K t12-tit-4 sub
\VUsaut{3.0mm}
\VUpk{10.2pt}{40}{ડોઇચ-આવ્રિકૂર સ્ટેશન.}
\VUsaut{3.0mm}

%%K t12-19-1 p
19. --- એક કારકુનના અવાજે મને એકાએક જગાડ્યો : « ડોઇચ-આવ્રિકૂર !
\textsuperscript{(*)} બધા ઊતરો ! » જકાતની તપાસ અને સરહદની બધી
કંટાળાજનક વિધિ થઈ. મારે મારું ખિસ્સાઘડિયાળ સ્ટેશનના \VUgras{મોટા
ઘડિયાળ} \textsuperscript{(89)} પ્રમાણે કરવું પડ્યું, જે પંચાવન મિનિટ
આગળ હતું ; માંડ જાગેલો હું \VUgras{મોટા કાંટાને} \textsuperscript{(90)}
\VUgras{નાનાથી} \textsuperscript{(91)} મુશ્કેલીથી જુદો પાડી શક્યો ;
\VUgras{આંકડાવાળી તકતી} \textsuperscript{(92)} પોતે સવારના ધુમ્મસને
લીધે તદ્દન ધૂંધળી હતી.

%%K t12-20-1 p
20. --- જકાતના અમલદારો તપાસ કરતા હતા ત્યારે હું બગાસાં ખાતાં ખાતાં
\VUgras{જાહેરાતનાં પોસ્ટર} \textsuperscript{(93)} ઉકેલતો હતો, જે
સ્ટેશનની દીવાલો શણગારે છે. સમય પસાર કરવા મેં નાસ્તો કર્યો, જર્મન
\VUgras{રેલજાળનો} \textsuperscript{(94)} નકશો જોયો કે સ્ટ્રાસબુર્ગ હજી
કેટલું દૂર છે, અને હું મારા ડબામાં પાછો ચડ્યો, મારી મુસાફરીનું ધ્યેય
જલદી પહોંચવાની આશાથી તાજો થઈને.

\VUsaut{6.0mm}
\VUfilet{22mm}
